\subsection{Spatial Summation and Classical Receptive Field Characterization}
Before investigating contextual modulations, we first validated the stability of the resting state and the classical receptive field (CRF) responses of the model. In the absence of targeted stimuli, the network maintained an asynchronous-irregular firing regime with stable baseline rates (e.g., $\sim 0.5$ Hz in superficial layer L2/3 and $\sim 6.2$ Hz in layer 5), confirming that the introduction of inter-columnar and inter-areal connectivity did not induce pathological oscillations. 

To characterize the CRF, we evaluated the network's response to variations in stimulus contrast, simulated by parametrically increasing the firing rate of the topographically aligned thalamic afferents (5, 10, 20, and 30 Hz). The model exhibited a monotonic, contrast-dependent scaling of firing rates, particularly prominent in the primary thalamorecipient layer (L4) and the superficial layer (L2/3). Furthermore, to ensure spatial specificity, a 20 Hz optimal stimulus was delivered exclusively to the central microcolumn (vertical preference). While the stimulated central column responded vigorously, the spatially distant peripheral column (horizontal preference) remained at baseline firing levels. This demonstrated that local horizontal connections effectively mediate lateral inhibition, preventing spurious horizontal propagation of strong sensory activation and preserving the spatial acuity of the CRF.

\subsection{Comparison of Architectural Configurations for Surround Suppression}
To systematically evaluate the structural substrates of extra-classical receptive field (eCRF) effects, we tested the network under a dual-stimulus contextual protocol. A central optimal stimulus (20 Hz) was presented at 500 ms, followed by an asynchronous peripheral stimulus of varying intensities (10, 20, and 30 Hz) at 1000 ms. 

We compared the magnitude of surround suppression across three distinct topological configurations. In the first configuration (V1 lateral connections only), the peripheral stimulus induced only a mild suppression. Under the highest contextual contrast (30 Hz), the L2/3 excitatory populations retained $63.3\% \pm 4.4\%$ of their peak CRF activity (equivalent to $\sim36.7\%$ suppression). Furthermore, layer 5 neurons showed virtually no suppressive modulation. This indicated that while intra-areal lateral connectivity contributes to contextual modulation, it is insufficient to replicate the high-magnitude suppression observed \textit{in vivo}.

The second configuration incorporated a single, global V2 module to test the hypothesis of non-specific top-down feedback. While this architecture successfully attenuated L5 activity—consistent with the anatomical targeting of feedback projections to infragranular layers—it failed to suppress L2/3. In fact, L2/3 activity remained at $84.0\% \pm 8.9\%$ of its CRF baseline. Analysis of the V2 module revealed that it acted as a linear integrator, indiscriminately pooling feedforward inputs from both the central and peripheral V1 columns without undergoing internal competition. Consequently, the feedback signal acted more as an excitatory gain amplifier rather than a selective suppressive filter.

\subsection{Robust Suppression via Hierarchical Competitive Dynamics}
The integration of a hierarchically segregated, competitive V2 architecture (Configuration 3) fundamentally altered the network's dynamics. In this design, two distinct V2 modules—each aligned to specific feature columns in V1—interacted via strong, reciprocal inhibitory connections. 

Under the identical eCRF stimulation protocol, the activation of the peripheral column triggered a robust winner-take-all dynamic in V2. The peripheral V2 module's activity surged, thereby actively suppressing the central V2 module, whose firing rate plummeted to near baseline levels ($\sim 1$ Hz). This highly selective refinement of the feedback signal profoundly modulated V1. In the central V1 microcolumn, L2/3 activity was suppressed to $22.9\% \pm 8.3\%$ of its maximum CRF response. This equates to a suppression index of approximately $77\%$, which quantitatively replicates the $\sim 75\%$ suppression magnitude historically recorded in macaque V1 (e.g., Busse et al., 2009). 

To conclusively isolate the source of this modulation, we performed virtual ablation experiments. Completely removing the horizontal lateral connections within V1 did not abolish the surround suppression effect. This remarkable resilience demonstrates that the top-down feedforward-feedback (FF-FB) loop, equipped with competitive inter-areal dynamics in higher visual cortices, is both necessary and sufficient to mediate robust far-surround suppression, overcoming the spatial and temporal limitations of slow, intra-areal lateral axons.

\subsection{Gamma-band Oscillations and E/I Balance}
Finally, we analyzed the temporal and oscillatory dynamics of the network during contextual modulation, leveraging a spatiotemporal kernel-based approximation to reconstruct the macroscopic Local Field Potential (LFP) from the unitary spike events (uLFP). 

Under normal physiological parameters, the emergence of robust surround suppression was accompanied by stable oscillatory activity in the high gamma-band range, peaking at approximately 75 Hz. To investigate the clinical applicability of the model, we simulated the excitation/inhibition (E/I) imbalance commonly associated with schizophrenia by systematically scaling down the probability of L2/3 inhibitory synapses (to 90\%, 75\%, and 50\% of their baseline strength). 

Perturbing this precise E/I balance not only caused a dramatic reduction in the depth of surround suppression but also triggered severe cortical hyperexcitability. Spectral analysis of the simulated LFP revealed that the disruption of local inhibition shifted the gamma oscillation peak down to $\sim 60$ Hz. At the severe 50\% inhibition threshold, the network lost its non-linear stability, exhibiting divergent low-frequency power surges during eCRF stimulation. These results provide a direct mechanistic link between microcircuit-level synaptic deficits, the degradation of contextual sensory filtering, and the macroscopic electrophysiological biomarkers characteristically observed in neuropsychiatric disorders.